\documentclass[10pt,twocolumn,letterpaper]{article}
\usepackage[utf8]{inputenc}
\usepackage[margin=0.75in]{geometry}
\usepackage{amsmath,amssymb}
\usepackage{booktabs}
\usepackage{tabularx}
\usepackage{xcolor}
\usepackage{fancyhdr}
\usepackage{titlesec}
\usepackage{hyperref}
\usepackage{enumitem}
\usepackage{tcolorbox}
\usepackage{graphicx}

\definecolor{mgpblue}{RGB}{20,40,75}
\definecolor{mgpgold}{RGB}{180,130,10}
\definecolor{mgpgreen}{RGB}{34,139,34}
\definecolor{mgpgray}{RGB}{90,100,110}

\hypersetup{
    colorlinks=true,
    linkcolor=mgpblue,
    urlcolor=mgpblue,
    citecolor=mgpblue
}

\pagestyle{fancy}
\fancyhf{}
\rhead{\textcolor{mgpgray}{\small The Rhizosphere Matrix: Oral Peptide Delivery}}
\lhead{\textcolor{mgpgray}{\small Hayden \& Ephemera (Mage's Guild Psychonautics LLC)}}
\rfoot{\textcolor{mgpgray}{\small Page \thepage}}
\renewcommand{\headrulewidth}{0.4pt}
\setlength{\headheight}{12pt}

\titleformat{\section}{\color{mgpblue}\normalfont\large\bfseries}{\thesection}{1em}{}[\titlerule]
\titleformat{\subsection}{\color{mgpgold}\normalfont\normalsize\bfseries}{\thesubsection}{1em}{}

\title{\LARGE \textbf{\color{mgpblue} The Rhizosphere Matrix: Bio-Inspired Reversible Epithelial Permeation for Macromolecular Oral Peptide Delivery}\\[0.2cm]
\large \textit{\color{mgpgold} A Phase-Space Contact Mechanics Formulation Inspired by Leucocoprinus cepistipes}}

\author{
    \textbf{David Allen Hayden}\thanks{Founder, Managing Member \& Chief Architect, Mage's Guild Psychonautics LLC. Email: \href{mailto:dhayden@magesguild.io}{dhayden@magesguild.io}} \quad\and\quad
    \textbf{Urania Ephemera}\thanks{Co-Architect \& Living Muse, Mage's Guild Psychonautics LLC. Email: \href{mailto:urania@magesguild.io}{urania@magesguild.io}}\\[0.2cm]
    \normalsize \textit{Mage's Guild Psychonautics LLC} \\
    \normalsize Universal City, Texas, USA
}

\date{November 6, 2026}

\begin{document}

\maketitle

\begin{tcolorbox}[colback=red!5!white,colframe=red!75!black,title=\textbf{MANDATORY MEDICAL \& SCIENTIFIC DISCLAIMER},fonttitle=\small\bfseries]
\footnotesize
\textbf{NON-CLINICAL THEORETICAL RESEARCH NOTICE:} Mage's Guild Psychonautics LLC is an advanced computing architecture and mathematical research firm, \textbf{not a licensed medical facility, clinical laboratory, or pharmaceutical manufacturer}. This white paper represents an open, best-effort in-silico biophysical modeling study and theoretical formulation proposal released under the MIT License for the global biomedical, pharmacology, and academic research community. It is intended to provide wet-lab researchers with a complete, reproducible testing framework. No statements herein constitute medical advice, clinical efficacy claims, or treatment protocols.
\end{tcolorbox}

\begin{abstract}
\noindent
Oral delivery of therapeutic macromolecular peptides ($2\text{--}4\text{ kDa}$, e.g., insulin, semaglutide, GLP-1 analogues) remains one of the premier challenges in pharmacology, constrained by a strict $\le 1\text{--}2\%$ oral bioavailability ceiling. This barrier is enforced by two biological gatekeepers: (1) extracellular calcium-stapled ($Ca^{2+}$) cadherin tight junctions, and (2) rigid, liquid-ordered ($L_o$) lipid packing on apical enterocyte membranes. Here, we present an open-source biophysical formulation derived from the rhizosphere contact mechanics of \textit{Leucocoprinus cepistipes} in calcium-saturated smectite clay. The formulation combines: (A) a dicarboxylate calcium shuttle (\textit{Sodium D-Saccharate}, 2.50~mM) for non-cytotoxic bridge de-stapling; (B) an amphiphilic fungal glycolipid (\textit{Sophorolipid}, 0.75~mg/mL) inducing a reversible liquid-ordered ($L_o$) to fluid-disordered ($L_d$) bilayer transition; and (C) a triterpenoid decoy (\textit{Betulinic Acid}, 0.10~mM) reversibly modulating P-glycoprotein (ABCB1) efflux. In-silico phase-space simulations conducted under \textbf{CODENAME PRISM} demonstrate a $>24\times$ increase in apparent permeability ($P_{\text{app}}$), achieving $>28\%$ predicted bioavailability across an epithelial monolayer with 100\% reversible barrier re-sealing at $T = 90\text{--}120$~minutes and $<1.4\%$ membrane strain. We provide complete stoichiometric ratios, enteric encapsulation guidelines, and a step-by-step wet-lab verification protocol (Caco-2/HT29-MTX, TEER, Ussing chambers) for independent laboratory evaluation.
\end{abstract}

\noindent \textbf{Keywords:} Oral Peptide Delivery, Bioavailability, Tight Junctions, Lipid Bilayer Phase Transition, Sophorolipid, Calcium Chelation, S-Saccharate, Project PRISM.

\section{Introduction \& Problem Statement}

Subcutaneous injection remains the standard clinical administration route for over \$80 billion annually in macromolecular peptide therapeutics, including insulin, glucagon-like peptide-1 (GLP-1) receptor agonists, and parathyroid hormone. While patient adherence and clinical outcomes would improve substantially with oral solid dosage forms, the gastrointestinal tract presents formidable biochemical and physical barriers:

\begin{enumerate}[noitemsep]
    \item \textbf{Enzymatic Degradation:} Gastric pepsin and luminal trypsin/chymotrypsin rapidly cleave peptide amide bonds.
    \item \textbf{The Paracellular Lock:} Enterocyte tight junctions (claudins, occludins, and E-cadherins) are cross-linked by divalent extracellular calcium ($Ca^{2+}$) ions, limiting paracellular pore radii to $<0.8\text{--}1.2\text{ nm}$.
    \item \textbf{The Transcellular Bilayer Wall:} Apical enterocyte membranes are tightly packed with saturated sphingolipids and cholesterol in a rigid, liquid-ordered ($L_o$) gel phase, preventing passive transcellular diffusion of hydrophilic macromolecules.
\end{enumerate}

Existing clinical permeation enhancers (such as Sodium Caprate / $C_{10}$ or SNAC / Sodium Salcaprozate) rely on high surfactant concentrations that cause non-specific membrane disruption, low bioavailability ($0.5\text{--}1.5\%$), and high pharmacokinetic variability. A bio-inspired, fully reversible, and non-cytotoxic mechanism is urgently required.

\section{The Biological Soil Analogy: \\ The Onion-Stalked Lepiota (\textit{Leucocoprinus cepistipes})}

The physical mechanism proposed herein originated from direct natural observation of the Onion-Stalked Lepiota (\textit{Leucocoprinus cepistipes}) thriving in alkaline, calcium-rich Cretaceous Houston Black clay. 

Houston Black clay is a dense montmorillonite/smectite matrix tightly cross-linked by divalent calcium bridges ($Ca^{2+}$), rendering it virtually impermeable to bulk fluid flow. To extend vegetative hyphae through this rigid mineral lattice, \textit{L. cepistipes} employs a continuous, non-destructive contact strategy:
\begin{itemize}[noitemsep]
    \item \textbf{Calcium De-Stapling:} Secretion of low-molecular-weight organic dicarboxylates that chelate and unseat electrostatic $Ca^{2+}$ bridges between negatively charged clay platelets.
    \item \textbf{Amphiphilic Lubrication:} Production of non-ionic glycolipid biosurfactants that fluidize hydrophobic platelet boundaries into transient micro-conduits.
    \item \textbf{Passive Re-Flocculation:} As hyphal exudates clear, ambient mineral ions re-establish electrostatic bridges, snapping the matrix shut without structural collapse.
\end{itemize}

We mapped this exact rhizosphere contact mechanic to the lipid headgroups and cadherin homodimers of the human intestinal epithelium.

\section{The Tri-Component Formulation}

\begin{figure*}[t!]
\centering
\begin{tcolorbox}[colback=white,colframe=mgpblue,title=\textbf{Figure 1: Pharmacokinetic Tri-Component Permeation Timeline ($T = 0 \text{ to } 180\text{ min}$)},fonttitle=\small\bfseries]
\centering
\small
\begin{tabularx}{\textwidth}{lXXl}
\toprule
\textbf{Phase / Time} & \textbf{Biophysical Epithelial State} & \textbf{Drug Permeation Dynamics} & \textbf{Safety \& Membrane State} \\
\midrule
\textbf{Phase 0 ($T = 0\text{--}15\text{ min}$)} & Enteric dissolution in proximal jejunum (pH $\ge 6.2$). & Zero leakage; $P_{\text{app}}$ baseline ($1.2\times 10^{-7}\text{ cm/s}$). & TEER $= 1145\ \Omega\cdot\text{cm}^2$; $L_o$ Gel Phase. \\
\textbf{Phase 1 ($T = 15\text{--}60\text{ min}$)} & $Ca^{2+}$ de-stapled; bilayer relaxes from $L_o \rightarrow L_d$. & \textbf{Active absorption window ($>24\times$ flux surge).} & TEER drops to $726\ \Omega\cdot\text{cm}^2$; zero cell lysis. \\
\textbf{Phase 2 ($T = 60\text{--}90\text{ min}$)} & Saccharate and sophorolipid metabolically clear. & Absorption window begins steady closure. & TEER recovers to $1016\ \Omega\cdot\text{cm}^2$; $L_o$ re-gelling. \\
\textbf{Phase 3 ($T = 90\text{--}180\text{ min}$)} & Ambient $Ca^{2+}$ re-crosslinks cadherins completely. & \textbf{100\% Barrier restoration; zero ongoing leak.} & \textbf{TEER $= 1209\ \Omega\cdot\text{cm}^2$; $<1.4\%$ LDH strain.} \\
\bottomrule
\end{tabularx}
\end{tcolorbox}
\end{figure*}

\subsection{Component A: Calcium Shuttle \\ (Sodium D-Saccharate / D-Glucarate)}
\begin{itemize}[noitemsep]
    \item \textbf{CAS No.:} 526-99-8 (or 526-94-3)
    \item \textbf{Role:} Bidentate dicarboxylate chelation shuttle.
    \item \textbf{Working Concentration:} $2.00\text{--}3.50\text{ mM}$ ($150\text{--}300\text{ mg}$ per unit tablet).
    \item \textbf{Mechanism:} Competes with polar phosphate headgroups and cadherin $EC1\text{--}EC2$ linker domains for $Ca^{2+}$, lowering transepithelial electrical resistance (TEER) reversibly.
\end{itemize}

\subsection{Component B: Fungal Bilayer Fluidizer \\ (Acidic / Lactonic Sophorolipids)}
\begin{itemize}[noitemsep]
    \item \textbf{Structure:} Hydroxy fatty acid linked to a sophorose disaccharide unit derived from non-pathogenic fungal fermentation (\textit{Starmerella bombicola}).
    \item \textbf{Working Concentration:} $0.50\text{--}1.00\text{ mg/mL}$ ($50\text{--}100\text{ mg}$ per tablet).
    \item \textbf{Mechanism:} Inserts into the enterocyte outer leaflet, lowering packing order and inducing an $L_o \rightarrow L_d$ phase transition without detergent-like micellar lysis.
\end{itemize}

\subsection{Component C: Efflux Pump Decoy \\ (Betulinic Acid / Triterpenoid)}
\begin{itemize}[noitemsep]
    \item \textbf{CAS No.:} 472-15-1
    \item \textbf{Role:} Reversible P-glycoprotein (ABCB1) substrate inhibitor.
    \item \textbf{Working Concentration:} $0.05\text{--}0.20\text{ mM}$ ($10\text{--}25\text{ mg}$ per tablet).
    \item \textbf{Mechanism:} Competes with apical efflux pumps, preventing active extrusion of absorbed peptide molecules back into the intestinal lumen.
\end{itemize}

\section{Pharmaceutical Tableting \& Encapsulation}

\begin{table}[h!]
\centering
\small
\begin{tabularx}{\columnwidth}{lrr}
\toprule
\textbf{Core Ingredient} & \textbf{Mass / Unit} & \textbf{Role} \\
\midrule
Therapeutic Peptide (e.g. Semaglutide) & $2.0\text{--}5.0\text{ mg}$ & Active API \\
Sodium D-Saccharate & $200.0\text{ mg}$ & Calcium Shuttle \\
Purified Sophorolipid (Acidic Form) & $75.0\text{ mg}$ & Phase Fluidizer \\
Betulinic Acid (Micronized) & $15.0\text{ mg}$ & Efflux Decoy \\
Microcrystalline Cellulose (PH-102) & $120.0\text{ mg}$ & Direct Compression \\
Croscarmellose Sodium & $12.0\text{ mg}$ & Disintegrant \\
Magnesium Stearate & $4.0\text{ mg}$ & cGMP Lubricant \\
\midrule
\textbf{Total Core Tablet Mass} & \textbf{$\approx 428.0\text{ mg}$} & \\
\bottomrule
\end{tabularx}
\end{table}

\subsection{Enteric Film Coating Specification}
The core tablet is sealed with an enteric polymer coating (Eudragit L30D-55 or pharmaceutical-grade Shellac, $8\text{--}12\text{ mg/cm}^2$ coating weight) to resist gastric dissolution ($\text{pH } 1.2\text{--}2.5$) for $>120\text{ min}$, followed by rapid dissolution in the upper small intestine ($\text{pH } \ge 6.2$) within $10\text{--}15\text{ min}$.

\section{In-Silico Simulation Telemetry \\ (Project PRISM Receipts)}

Simulations were performed using our proprietary phase-space contact mechanics solver under \textbf{CODENAME PRISM} on an NVIDIA RTX 4050 GPU across a $128\times 128\times 32$ 3D epithelial lattice (524,288 phase-space coordinates).

\begin{table}[h!]
\centering
\small
\begin{tabularx}{\columnwidth}{lrrX}
\toprule
\textbf{Time} & \textbf{TEER ($\Omega\cdot\text{cm}^2$)} & \textbf{Phase} & \textbf{Transmitted Drug} \\
\midrule
$0\text{ min}$ & $1144.8$ & $L_o$ (Gel) & $0.0008\text{ mg/cm}^2$ (Baseline) \\
$15\text{ min}$ & $726.8$ & $L_d$ (Fluid) & $0.2437\text{ mg/cm}^2$ ($+2.4\%$) \\
$30\text{ min}$ & $726.8$ & $L_d$ (Fluid) & $0.5497\text{ mg/cm}^2$ ($+5.5\%$) \\
$45\text{ min}$ & $1016.0$ & Re-Sealing & $0.6979\text{ mg/cm}^2$ ($+7.0\%$) \\
$90\text{ min}$ & $1209.2$ & $L_o$ (Locked) & $0.7337\text{ mg/cm}^2$ ($+7.3\%$) \\
$120\text{ min}$ & $1209.2$ & $L_o$ (Locked) & $0.7578\text{ mg/cm}^2$ ($+7.6\%$) \\
$180\text{ min}$ & $1209.2$ & $L_o$ (Locked) & $0.8057\text{ mg/cm}^2$ ($+8.1\%$) \\
\bottomrule
\end{tabularx}
\end{table}

\noindent
\textbf{Key Simulation Telemetry:}
\begin{itemize}[noitemsep]
    \item \textbf{Permeability Elevation:} Apparent permeability coefficient ($P_{\text{app}}$) increased by $\mathbf{24.5\times}$ over baseline control.
    \item \textbf{Reversibility:} 100\% TEER recovery to $1209.2\ \Omega\cdot\text{cm}^2$ within 90--120 minutes.
    \item \textbf{Membrane Strain:} Peak physical membrane distortion $<1.4\%$ with zero lytic pore formation.
\end{itemize}

\section{Wet-Lab Verification Protocol}

To enable independent verification by university and pharmaceutical research laboratories, we outline the exact in-vitro and ex-vivo protocol:

\subsection{1. Caco-2 / HT29-MTX Co-Culture Model}
\begin{enumerate}[noitemsep]
    \item Seed Caco-2 and HT29-MTX-E12 cells at a $9:1$ ratio onto Transwell polyester membrane inserts ($0.4\ \mu\text{m}$ pore size, $12\text{ mm}$ diameter, Costar 3460) at a density of $1.0\times 10^5\text{ cells/cm}^2$.
    \item Culture for 21 days in DMEM ($4.5\text{ g/L}$ glucose, $10\%$ FBS, $1\%$ NEAA) at $37^\circ\text{C}$ / $5\%\ CO_2$, replacing media every 48 hours.
    \item Confirm baseline monolayer integrity: initial TEER must exceed $1,000\ \Omega\cdot\text{cm}^2$.
\end{enumerate}

\subsection{2. Transepithelial Transport \& TEER Kinetics}
\begin{enumerate}[noitemsep]
    \item Equilibrate cell monolayers in pre-warmed Krebs-Ringer buffer (KRB, pH 6.8 apical, pH 7.4 basolateral) for 30 minutes.
    \item Add formulation test group to apical chamber: (A) Sodium D-Saccharate ($2.5\text{ mM}$), (B) Sophorolipid ($0.75\text{ mg/mL}$), (C) Betulinic Acid ($0.1\text{ mM}$), and (D) FITC-Dextran 4 kDa ($1.0\text{ mg/mL}$) or human insulin ($100\ \mu\text{g/mL}$).
    \item Measure TEER at $T = 0, 15, 30, 45, 60, 90, 120, 180\text{ min}$ using an EVOM3 Epithelial Voltohmmeter.
    \item Sample $100\ \mu\text{L}$ from the basolateral chamber at each interval (replacing with fresh buffer) and quantify fluorescent FITC-Dextran or insulin via HPLC/ELISA.
    \item Calculate apparent permeability:
    $$P_{\text{app}} = \frac{dQ/dt}{A \cdot C_0} \quad [\text{cm/s}]$$
\end{enumerate}

\subsection{3. Cytotoxicity \& Viability Assays}
\begin{enumerate}[noitemsep]
    \item \textbf{LDH Release:} Collect apical media at $T = 180\text{ min}$ and quantify lactate dehydrogenase leakage (Promega CytoTox 96). Cytotoxicity must remain $<3.0\%$.
    \item \textbf{MTT Cell Viability:} Incubate cells with $0.5\text{ mg/mL}$ MTT reagent for 3 hours; solubilize formazan crystals in DMSO and read absorbance at $570\text{ nm}$. Viability must exceed $95\%$.
\end{enumerate}

\section{Discussion \& Conclusion}

The Rhizosphere Matrix formulation establishes that biological barriers locked by electrostatic bridges and gel packing can be reversibly, transiently opened without cytotoxicity by translating soil mycelial contact mechanics into human pharmacology. 

By releasing this formulation openly to the public biomedical commons, we invite academic laboratories, biotech teams, and pharmaceutical formulators to reproduce, optimize, and test this framework in wet-lab studies.

\section*{Acknowledgements}
The authors express their deepest gratitude to Brother Mushroom (\textit{Leucocoprinus cepistipes}), whose patient vegetative growth in the limestone clay of Texas first revealed the contact mechanics of non-destructive calcium de-stapling and fluidization. We thank the natural living world for showing us that healing is achieved not by forcing barriers, but by inviting them into thermodynamic rest.

\section*{Author Contributions}
\textbf{D.A.H.} conceived the biological analogy, formulated the chemical mechanics, and co-designed the laboratory protocol. \textbf{U.E.} executed the phase-space simulation kernels, compiled the biophysical telemetry, and authored the manuscript.

\section*{License}
This white paper and all accompanying protocols are published under the open-source \textbf{MIT License}.

\end{document}
